\section{Abstract}
\label{sec:Abstract}

GRACE steckt in der Pilotphase. Das bringt Probleme mit sich: Formale Dokumentation existiert kaum, Konfigurationswissen ist über Azure DevOps Wikis, E-Mails, Git-Historien und zehn Jahre persönliche Erfahrung verstreut. Die Folgen zeigen sich deutlich: Die Vollständigkeit liegt bei nur 60-70\%, durchschnittlich treten 10-15 Fehler pro Konfiguration auf, die Durchlaufzeit beträgt rund acht Wochen. Das implizite Wissen liegt in den Köpfen weniger Experten und droht bei Personalwechseln verloren zu gehen \citep{polanyi:tacit_dimension1966, nonaka:knowledge_creating_company1995}.
\\\\
\textbf{Wäre es möglich, diesen Konfigurationsprozess systematisch zu verbessern?} Diese Frage stand am Anfang der Arbeit. Das Ziel: Eine Lösung entwickeln, die fragmentiertes Expertenwissen externalisiert und für LLM-Automatisierung nutzbar macht. Das Ergebnis ist eine Drei-Artefakte-Architektur, entwickelt nach Action Research \citep{lewin:action_research1946}. Sie besteht aus drei Komponenten: Dem GRACE Onboarding Fragebogen (aufgebaut nach \cite{devellis:scale_development2017}), dem Entity Mapping Workbook (basierend auf \cite{chen:erm1976}) und den Standard Operating Procedures. Auf dieser Wissensbasis arbeitet der Master Creator, ein konversationelles Tool mit lokalem 7B-Modell, das Datensouveränität gewährleistet \citep{hummel:data_sovereignty2021, gupta:multilingual2025}.
\\\\
Die Auswertung der Ergebnisse zeigt beachtliche Verbesserungen. Ein kontrolliertes Experiment mit 18 Datenpunkten ergab eine Zeitreduktion von 47\%, die Fehlerrate sank um 81\%. Die Vollständigkeit stieg von anfänglich 60-70\% auf 98,1\%. Besonders aufschlussreich ist die zweistufige Analyse: Bereits die reine Strukturierung des Wissens verbesserte die Vollständigkeit auf 86,1\%. Dies bildete eine notwendige Basis. Die anschließende LLM-Automatisierung steigerte den Wert dann auf die finalen 98,1\%. Das verwendete lokale 7B-Modell erreichte trotz vergleichsweise bescheidener Hardware-Anforderungen durchgängig referenzielle Integrität über alle sechs Entitätstypen hinweg.
\\\\
Die Arbeit zeigt: LLM-gestützte Konfigurationsansätze funktionieren für APS-Systeme unter realen Produktionsbedingungen. Die zentrale Erkenntnis dabei ist, dass Wissensexternalisierung keine optionale Ergänzung darstellt. Sie ist notwendige Voraussetzung für erfolgreiche LLM-Automatisierung in komplexen Enterprise-Domänen.

